% Part: many-valued-logic
% Chapter: syntax-and-semantics
% Section: connectives

\documentclass[../../../include/open-logic-section]{subfiles}

\begin{document}

\olfileid{mvl}{syn}{con}

\olsection{اللغات والروابط}

يستعمل منطق القضايا الكلاسيكي، ومنطقات كثيرة أخرى، مخزونًا محددًا من
\emph{الثوابت القضايا} و\emph{الروابط}. فنستعمل مثلًا ما يأتي بوصفه
رموزًا أولية:
\begin{enumerate}
\tagitem{prvFalse}{الثابت القضوي لـ!!{falsity}~$\lfalse$.}{}
\tagitem{prvTrue}{الثابت القضوي لـ!!{truth}~$\ltrue$.}{}
\item الروابط المنطقية:
  \startycommalist
  \iftag{prvNot}{\ycomma $\lnot$ (النفي)}{}%
  \iftag{prvAnd}{\ycomma $\land$ (الاقتران)}{}%
  \iftag{prvOr}{\ycomma $\lor$ (الفصل)}{}%
  \iftag{prvIf}{\ycomma $\lif$ (!!{conditional})}{}%
  \iftag{prvIff}{\ycomma $\liff$ (!!{biconditional})}{}%
\end{enumerate}
\iftag{defNot,defOr,defAnd,defIf,defIff,defTrue,defFalse,defEx,defAll}{%
إضافة إلى الروابط الأولية أعلاه، نستعمل أيضًا رموزًا معرّفة بوصفها
اختصارات، مثل
\startycommalist
  \iftag{defNot}{\ycomma $\lnot$ (النفي)}{}%
  \iftag{defAnd}{\ycomma $\land$ (الاقتران)}{}%
  \iftag{defOr}{\ycomma $\lor$ (الفصل)}{}%
  \iftag{defIf}{\ycomma $\lif$ (!!{conditional})}{}%
  \iftag{defIff}{\ycomma $\liff$ (!!{biconditional})}{}%
  \iftag{defFalse}{\ycomma $\lfalse$ (!!{falsity})}{}%
  \iftag{defTrue}{\ycomma $\ltrue$ (!!{truth})}.}{}

تستعمل الروابط نفسها في المنطقات المتعددة القيم أيضًا. غير أنه من
المفيد في كثير من الأحيان إدراج صيغ مختلفة لرابط، كالاقتران مثلًا،
في المنطق نفسه، وهذا يتطلب رموزًا مختلفة لتمييز الصيغ بعضها من بعض.
وتتضمن بعض المنطقات المتعددة القيم كذلك روابط لا مقابل لها في المنطق
الكلاسيكي. لذلك سنكون أعم قليلًا من المعتاد.

\begin{defn}
  تتألف \emph{لغة قضايا} من مجموعة~$\Lang L$ من \emph{الروابط}.
  ولكل رابط~$\star$ \emph{رتبة}؛ ويقال عن رابط رتبته~$n$ إنه
  \emph{ذو~$n$ من المواضع}. وتسمى الروابط ذات الرتبة~$0$
  \emph{ثوابت} أيضًا؛ وتسمى الروابط ذات الرتبة~$1$ \emph{أحادية}،
  والروابط ذات الرتبة~$2$ \emph{ثنائية}.
\end{defn}

\begin{ex}
  تتألف اللغة القياسية لمنطق القضايا~$\Lang L_0$ من الروابط الآتية
  (مع رتبها):
  $\lfalse$~($0$)،
  $\lnot$~($1$)،
  $\land$~($2$)،
  $\lor$~($2$)،
  $\lif$~($2$). وستستعمل معظم المنطقات التي ننظر فيها هذه اللغة.
  وتستعمل بعض المنطقات، بحكم العرف والاصطلاح، رموزًا مختلفة لبعض
  الروابط. ففي منطق الضرب مثلًا يكون رمز الاقتران غالبًا~$\odot$ بدل
  $\land$. ومن الملائم أحيانًا إضافة مؤثر جديد، مثل مؤثر التعيّن
  $\triangle$ (ذو موضع واحد).
\end{ex}

\end{document}
